% Wersja źródłowa: audyt II, jednostka 32
\setcounter{section}{31}
\section{Teoria spektralna operatorów ograniczonych na przestrzeniach Hilberta}
\label{sec:spectral-hilbert32}
\noindent\textbf{Wyniki wykorzystywane.}
Korzystamy z przemiennego twierdzenia Gelfanda--Naimarka
(Twierdzenie~\ref{thm:commutative-cstar31}), twierdzenia
Stone'a--Weierstrassa (Lemat~\ref{lem:stone-weierstrass31}),
twierdzenia Riesza--Markowa o reprezentacji dodatnich
funkcjonałów na $C(K)$, własności całki Lebesgue'a
i twierdzenia Riesza--Schaudera (Twierdzenie~\ref{thm:riesz-schauder30}).
\par\medskip

Poprzedni rozdział dostarczył holomorficznego rachunku funkcyjnego
w dowolnej zespolonej algebrze Banacha i twierdzenia Gelfanda--Naimarka
dla przemiennych algebr $C^*$. Dla operatorów normalnych na przestrzeniach
Hilberta otrzymamy znacznie więcej: rachunek wszystkich ograniczonych
funkcji borelowskich oraz reprezentację przez \emph{miarę spektralną}.
Nie zakładamy ośrodkowości przestrzeni Hilberta. Iloczyn skalarny jest
liniowy w pierwszym argumencie, tak jak w \hyperref[sec:hilbert13]{rozdziale~13}.

\subsection{Operatory normalne i ciągły rachunek funkcyjny}

\begin{definicja}[Operator normalny, samosprzężony i unitarny]
\label{def:normal32}
Niech $H$ będzie niezerową zespoloną przestrzenią Hilberta,
a $T\in\mathcal L(H,H)$. Operator $T$ jest \emph{normalny}, gdy
$TT^*=T^*T$, \emph{samosprzężony}, gdy $T=T^*$, oraz
\emph{unitarny}, gdy $T^*T=TT^*=I_H$. Operator samosprzężony
$T$ jest \emph{dodatni}, co zapisujemy $T\geq0$, gdy
$\langle Tx,x\rangle\geq0$ dla każdego $x\in H$.
\end{definicja}

\begin{lemat}[Zgodność widma z podalgebrą $C^*$]
\label{lem:spectral-permanence32}
Niech $H$ będzie niezerową zespoloną przestrzenią Hilberta,
$T\in\mathcal L(H,H)$ operatorem normalnym i
$A=C^*(I_H,T)$ domkniętą w normie podalgebrą generowaną
przez $I_H,T,T^*$. Wówczas $A$ jest przemienną unitalną
algebrą $C^*$ i $\sigma_A(T)=\sigma_{\mathcal L(H,H)}(T)$.
\end{lemat}
\begin{proof}
Normalność daje przemienność $T$ i $T^*$, a więc całej
podalgebry $A$. Zawsze $\sigma_{\mathcal L(H,H)}(T)\subseteq\sigma_A(T)$.
Niech odwrotnie $\lambda\notin\sigma_{\mathcal L(H,H)}(T)$ i
$B=T-\lambda I_H$. Wówczas $\|x\|\le\|B^{-1}\|\|Bx\|$, więc dla
$m=\|B^{-1}\|^{-2}>0$ zachodzi
$\|Bx\|^2\geq m\|x\|^2$ dla wszystkich $x\in H$.
Operator $S=B^*B$ należy do $A$, jest dodatni i spełnia
$mI_H\leq S\leq\|S\|I_H$. Przypomnijmy uzasadnienie wzoru na normę operatora
samosprzężonego $R$. Niech
$c=\sup_{\|x\|=1}|\langle Rx,x\rangle|$.
Dla jednostkowych $x,y$ dobieramy fazę wektora $y$ tak,
aby $\langle Rx,y\rangle$ było nieujemne rzeczywiste.
Z tożsamości
\[
4\operatorname{Re}\langle Rx,y\rangle
=\langle R(x+y),x+y\rangle-\langle R(x-y),x-y\rangle
\]
i prawa równoległoboku dostajemy
$4|\langle Rx,y\rangle|\le c(\|x+y\|^2+\|x-y\|^2)=4c$.
Supremum po $x,y$ daje $\|R\|\le c$, a odwrotna
nierówność wynika z nierówności Cauchy'ego--Schwarza.
Stosując ten wzór do $R=I_H-S/\|S\|$, otrzymujemy
$\|I_H-S/\|S\|\|\leq1-m/\|S\|<1$.
Szereg Neumanna dowodzi, że $S^{-1}\in A$. Ponieważ $B$ jest
normalny, $B$ i $B^*$ komutują z $S^{-1}$; zatem
$S^{-1}B^*$ jest obustronną odwrotnością $B$ w $A$.
W konsekwencji $\lambda\notin\sigma_A(T)$.
\end{proof}

\begin{twierdzenie}[Ciągły rachunek funkcyjny operatora normalnego]
\label{thm:continuous-calculus32}
Niech $H$ i $T$ będą jak w lemacie~\ref{lem:spectral-permanence32},
a $K=\sigma(T)\subset\C$. Istnieje dokładnie jeden izometryczny
unitalny homomorfizm algebr z inwolucją
$\Phi_T:C(K,\C)\to\mathcal L(H,H)$ taki, że
$\Phi_T(\mathrm{id}_K)=T$. Dla $f\in C(K,\C)$ piszemy $f(T)=\Phi_T(f)$.
Zachodzą
$\|f(T)\|=\max_{z\in K}|f(z)|$,
$\overline f(T)=f(T)^*$ i $\sigma(f(T))=f(K)$.
Jeżeli $p(z,\overline z)$ jest wielomianem w dwóch zmiennych,
to $p(T, T^*)=p|_K(T)$.
\end{twierdzenie}
\begin{proof}
Przemienne twierdzenie Gelfanda--Naimarka z
poprzedniego rozdziału daje izometryczny
izomorfizm $A\simeq C(\Delta(A),\C)$, zaś lemat
\ref{lem:spectral-permanence32} identyfikuje $K$ z widmem $T$
w $A$. Odwzorowanie
$\Delta(A)\ni\chi\mapsto\chi(T)\in K$ jest ciągłą surjekcją.
Jest różnowartościowe: charakter jest określony przez wartości
na $T,T^*$ i jednostce, a $\chi(T^*)=\overline{\chi(T)}$.
Jest więc homeomorfizmem zwartych przestrzeni Hausdorffa.
Przeniesienie transformacji Gelfanda daje $\Phi_T$.
Jednoznaczność wynika z gęstości wielomianów w $z,\overline z$
w $C(K)$, czyli z twierdzenia Stone'a--Weierstrassa.
Izometria i zgodność z inwolucją wynikają z twierdzenia
Gelfanda--Naimarka. Dla równości widm rozważmy
$B=C^*(I_H,f(T))\subseteq A$. Operator $f(T)$ jest
normalny. Lemat~\ref{lem:spectral-permanence32} dla
$f(T)$ daje $\sigma_{\mathcal L(H,H)}(f(T))=\sigma_B(f(T))$,
a inkluzje widm przy włączaniu algebr dają
\[
\sigma_{\mathcal L(H,H)}(f(T))\subseteq\sigma_A(f(T))
\subseteq\sigma_B(f(T)).
\]
Wszystkie trzy zbiory są więc równe. W $A\simeq C(K)$
widmo $f(T)$ jest dokładnie $f(K)$.
\end{proof}

\begin{uwaga}[Związek z rachunkiem holomorficznym]
Jeżeli $f$ jest holomorficzna na otoczeniu $K$, to $f(T)$
z twierdzenia~\ref{thm:continuous-calculus32} zgadza się z całką
Dunforda z poprzedniego rozdziału.
Istotnie, dla $w\notin K$ rachunek ciągły daje
$(w-z)^{-1}(T)=(wI_H-T)^{-1}$. Jeżeli $\Gamma$ jest
konturem dopuszczalnym w definicji całki Dunforda,
ciągłość homomorfizmu $\Phi_T$ pozwala przenieść go
pod znak całki (jako granicę sum Riemanna). Wtedy
\[
 \frac1{2\pi i}\int_\Gamma f(w)(wI_H-T)^{-1}\,\mathrm dw
 =\Phi_T\!\left(\frac1{2\pi i}\int_\Gamma
 \frac{f(w)}{w-z}\,\mathrm dw\right)=\Phi_T(f|_K),
\]
gdzie ostatnia równość jest skalarnym wzorem Cauchy'ego
stosowanym punktowo dla $z\in K$.
Dla operatora nienormalnego sam rachunek
ciągły na $C(\sigma(T))$ z własnością $\overline f(T)=f(T)^*$
na ogół nie istnieje.
\end{uwaga}

\subsection{Miary spektralne i twierdzenie spektralne}

\begin{definicja}[Miara spektralna]
\label{def:pvm32}
Niech $H$ będzie zespoloną przestrzenią Hilberta i $K\subset\C$
niepustym zwartym zbiorem. \emph{Miarą spektralną} na
$\mathcal B(K)$ nazywamy odwzorowanie
$E:\mathcal B(K)\to\mathcal L(H,H)$ spełniające:
$E(B)$ jest rzutem ortogonalnym dla każdego zbioru borelowskiego $B$,
$E(\varnothing)=0$, $E(K)=I_H$ oraz
$E(B\cap C)=E(B)E(C)$ dla $B,C\in\mathcal B(K)$.
Ponadto dla parami rozłącznych $(B_n)_{n\geq1}$ wymaga się
\emph{ciągowej addytywności w silnej topologii operatorowej}, tj.
$E(\bigcup_nB_n)x=\lim_{N\to\infty}\sum_{n=1}^NE(B_n)x$
dla każdego $x\in H$.
Dla $x\in H$ związana z $E$ dodatnia skończona miara skalarna
to $\mu_x(B):=\langle E(B)x,x\rangle=\|E(B)x\|^2$.
\end{definicja}

\begin{lemat}[Model cykliczny]
\label{lem:cyclic-model32}
Niech $T\in\mathcal L(H,H)$ będzie normalny, $K=\sigma(T)$,
$a\in H\setminus\{0\}$ i
$H_a=\overline{\{f(T)a:f\in C(K,\C)\}}$.
Wówczas $H_a$ jest podprzestrzenią redukującą $T$
(niezmienniczą względem $T$ i $T^*$).
Istnieje jedyna skończona dodatnia regularna miara borelowska
$\nu_a$ na $K$ i izometria liniowa na
$U_a:L^2(K,\nu_a)\to H_a$ takie, że
$U_a1=a$ i $U_aM_z=(T|_{H_a})U_a$, gdzie
$(M_zf)(w)=wf(w)$.
\end{lemat}
\begin{proof}
Z domknięcia algebry $C(K)$ na sprzężenie otrzymujemy
$T H_a\subseteq H_a$ i $T^*H_a\subseteq H_a$.
Definiujemy $L_a(f)=\langle f(T)a,a\rangle$. Jest to ciągły
dodatni funkcjonał na $C(K)$, ponieważ dla $f\geq0$
możemy zapisać $f=|\sqrt f|^2$, a więc
$f(T)=(\sqrt f(T))^*\sqrt f(T)\geq0$.
Z twierdzenia Riesza--Markowa o reprezentacji dodatnich
funkcjonałów na $C(K)$ otrzymujemy jednoznaczną dodatnią,
skończoną regularną miarę borelowską $\nu_a$ spełniającą
$L_a(f)=\int_K f(z)\,\mathrm d\nu_a(z)$.
Mamy $\nu_a(K)=\|a\|^2$ oraz
\[
 \langle f(T)a,g(T)a\rangle
 =\langle (\overline g f)(T)a,a\rangle
 =\int_K f(z)\overline{g(z)}\,\mathrm d\nu_a(z).
\]
Zatem odwzorowanie $f\mapsto f(T)a$ zależy tylko od
klasy $f$ w $L^2(\nu_a)$ i jest izometrią liniową na
swój obraz w $H$. Nie zakładamy, że półnorma $L^2$
jest normą na wszystkich funkcjach ciągłych.
Aby wykazać gęstość ich klas, dla borelowskiego $B$
i $\varepsilon>0$ wybieramy z regularności miary
zwarty $C\subseteq B$ i względnie otwarty $O\supseteq B$
z $\nu_a(O\setminus C)<\varepsilon^2$.
Lemat Urysohna daje $h\in C(K,[0,1])$ równe $1$ na $C$
i $0$ poza $O$, więc
$\|h-\chi_B\|_{L^2(\nu_a)}^2\le\nu_a(O\setminus C)<\varepsilon^2$.
Gęstość funkcji prostych w $L^2$ kończy argument. Izometria przedłuża się do $U_a$ na całym
$L^2(K,\nu_a)$; jej obraz jest domknięty i równy $H_a$.
Ponieważ $U_a(zf)=T U_a(f)$ dla funkcji ciągłych,
przez gęstość otrzymujemy równość operatorową.
Jeżeli inna miara i izometria na spełniają te warunki,
to izometria jest unitarna, zatem przeplata również
sprzężone operatory mnożenia. Dla wielomianu
$p(z,\overline z)$ obraz funkcji $p$ musi być równy
$p(T,T^*)a$. Gęstość tych wielomianów w $C(K)$
i skończoność miar dają to samo dla każdej funkcji ciągłej.
Parowanie z $a$ identyfikuje całki funkcji ciągłych
z $L_a$, więc jednoznaczność Riesza--Markowa daje
równość miar. Gęstość w $L^2$ daje równość izometrii.
\end{proof}

\begin{lemat}[Całkowanie względem dowolnej miary spektralnej]
\label{lem:pvm-integral32}
Niech $E$ będzie miarą spektralną z Definicji~\ref{def:pvm32}.
Dla każdej ograniczonej borelowskiej funkcji $f$ całka
$\int_K f\,\mathrm dE$ jest dobrze określona jako granica
w normie operatorowej całek funkcji prostych aproksymujących
$f$ jednostajnie. Jest to unitalny homomorfizm algebr
z inwolucją, o normie nie większej niż $1$.
Dla $\mu_{x,y}(B)=\langle E(B)x,y\rangle$ zachodzą
\[
\left\langle\left(\int_K f\,\mathrm dE\right)x,y\right\rangle
=\int_K f\,\mathrm d\mu_{x,y},\qquad
\left\|\left(\int_K f\,\mathrm dE\right)x\right\|^2
=\int_K |f|^2\,\mathrm d\mu_x .
\]
\end{lemat}
\begin{proof}
Dla $s=\sum_{j=1}^m c_j\chi_{B_j}$ na skończonym
borelowskim podziale $K$ definiujemy
$I_E(s)=\sum_{j=1}^m c_j E(B_j)$. Skończona addytywność
i wspólne rozdrobnienie podziałów dają niezależność
od zapisu, liniowość i multiplikatywność; samosprzężoność
rzutów daje $I_E(\overline s)=I_E(s)^*$.
Ortogonalność obrazów rzutów daje
\[
\|I_E(s)x\|^2=\sum_j |c_j|^2\|E(B_j)x\|^2
\le\|s\|_\infty^2\|x\|^2 .
\]
Ograniczoną funkcję borelowską aproksymujemy jednostajnie
funkcjami prostymi, dzieląc ograniczony zbiór jej wartości
na skończenie wiele borelowskich zbiorów o średnicy
dążącej do zera. Oszacowanie normy i zupełność
$\mathcal L(H,H)$ dają granicę niezależną od aproksymacji.
Ciągłość mnożenia i sprzężenia przenosi własności na granicę.
Silna addytywność $E$ daje ciągową addytywność
$\mu_{x,y}$. Dla skończonego rozłącznego układu $(B_j)$
\[
\sum_j|\langle E(B_j)x,y\rangle|
=\sum_j|\langle E(B_j)x,E(B_j)y\rangle|
\le\|x\|\|y\|,
\]
na mocy nierówności Cauchy'ego--Schwarza dla sum i
ortogonalności. Miara zespolona ma zatem skończone
wahanie całkowite. Dla funkcji prostych oba wzory
wynikają z definicji, a jednostajna aproksymacja
i skończoność miar dają oba wzory dla $f$.
\end{proof}

\begin{twierdzenie}[Twierdzenie spektralne dla operatora normalnego]
\label{thm:spectral-normal32}
Niech $H$ będzie niezerową zespoloną przestrzenią Hilberta,
$T\in\mathcal L(H,H)$ operatorem normalnym, a $K=\sigma(T)$.
Istnieje dokładnie jedna miara spektralna
$E:\mathcal B(K)\to\mathcal L(H,H)$ taka, że
$T=\int_Kz\,\mathrm dE(z)$.
Dla każdej ograniczonej borelowskiej funkcji $f:K\to\C$
istnieje operator $f(T)=\int_K f(z)\,\mathrm dE(z)$ spełniający
$\langle f(T)x,y\rangle=\int_K f(z)\,\mathrm d\mu_{x,y}(z)$,
gdzie $\mu_{x,y}(B):=\langle E(B)x,y\rangle$
jest skończoną miarą zespoloną, otrzymaną przez
polaryzację miar dodatnich $\mu_x$.
Zachodzą: $(fg)(T)=f(T)g(T)$,
$f(T)^*=\overline f(T)$ oraz
$\|f(T)x\|^2=\int_K|f(z)|^2\,\mathrm d\mu_x(z)$.
Rachunek ten przedłuża ciągły rachunek funkcyjny z
twierdzenia~\ref{thm:continuous-calculus32}.
\end{twierdzenie}
\begin{proof}
\textbf{Krok 1: rozkład na podprzestrzenie cykliczne.}
Rozważamy rodziny parami ortogonalnych, niezerowych podprzestrzeni
postaci $H_a$ z lematu~\ref{lem:cyclic-model32}. Są one redukujące.
Porządkujemy te rodziny przez włączanie; suma łańcucha
jest rodziną parami ortogonalną. Lemat Zorna daje
rodzinę maksymalną $(H_{a_j})_{j\in J}$.
Ortogonalne dopełnienie domkniętej sumy jest redukujące,
bo suma jest niezmiennicza względem $T$ i $T^*$.
Każdy niezerowy wektor z dopełnienia generowałby kolejną
podprzestrzeń cykliczną: dopełnienie jest niezmiennicze
względem wielomianów w $T,T^*$, a przez aproksymację
także względem każdego $f(T)$ dla $f\in C(K)$.
Przeczy to maksymalności.
Zatem $H=\bigoplus_{j\in J}H_{a_j}$, także dla nieprzeliczalnego $J$.

\textbf{Krok 2: konstrukcja miary.}
Z lematu~\ref{lem:cyclic-model32} każdy $T|_{H_{a_j}}$
jest unitarnie równoważny operatorowi mnożenia $M_z$ na
$L^2(K,\nu_{a_j})$. Oznaczmy tę równoważność przez $U_j$.
Dla borelowskiego $B\subseteq K$ definiujemy
$E(B)|_{H_{a_j}}=U_jM_{\chi_B}U_j^{-1}$.
Są to rzuty ortogonalne, komutują i spełniają
$E(B\cap C)=E(B)E(C)$, $E(K)=I_H$. Jeżeli $(B_n)$ jest
rodziną parami rozłączną, to dla każdego $x_j\in H_{a_j}$
\[
 \bigl\|\bigl(E(\textstyle\bigcup_n B_n)
 -\sum_{n=1}^{N}E(B_n)\bigr)x_j\bigr\|^2
 =\int_{\bigcup_{n>N}B_n}|U_j^{-1}x_j|^2\,\mathrm d\nu_{a_j}
 \longrightarrow0.
\]
Każdy wektor $x\in\bigoplus_{j\in J}H_{a_j}$ ma co najwyżej
przeliczalnie wiele niezerowych składowych; sumując powyższe
oszacowania i stosując zbieżność zmajoryzowaną dla szeregu
nieujemnych składników (każdy jest ograniczony przez
$\|x_j\|^2$, a $\sum_j\|x_j\|^2=\|x\|^2$),
otrzymujemy silną ciągową addytywność.

\textbf{Krok 3: rachunek funkcyjny.}
Dla ograniczonej funkcji borelowskiej $f$ definiujemy
$f(T)|_{H_{a_j}}=U_jM_fU_j^{-1}$.
Normy tych operatorów nie przekraczają $\|f\|_\infty$,
więc określają jeden operator na całej sumie ortogonalnej.
Dla funkcji prostych tak określony operator jest równy
$I_E(f)$ z Lematu~\ref{lem:pvm-integral32}. Dla
ograniczonych funkcji borelowskich jednostajna aproksymacja
i oszacowanie normy operatorów mnożenia dają tę samą
równość. Wszystkie własności całkowe wynikają z tego lematu.
Dla $f(z)=z$ dostajemy $T$. Dla funkcji ciągłych
otrzymujemy wcześniejszy rachunek, gdyż oba homomorfizmy
zgadzają się na $z$ i $\overline z$, a te generują
$C(K)$ jako zamkniętą algebrę.

\textbf{Krok 4: jednoznaczność.}
Jeżeli $F$ jest inną miarą spektralną o podanej własności,
to niezależnie zdefiniowany rachunek z
Lematu~\ref{lem:pvm-integral32} oparty na $F$
przyporządkowuje wielomianom
w $z,\overline z$ operatory $p(T,T^*)$.
Z twierdzenia Stone'a--Weierstrassa to samo zachodzi
dla wszystkich $f\in C(K)$. Dla dowolnych $x,y\in H$
miary zespolone $\langle E(\cdot)x,y\rangle$ i
$\langle F(\cdot)x,y\rangle$ mają więc identyczne całki
z wszystkich funkcji ciągłych. Każda z nich jest regularna:
wynika to przez polaryzację ze skończoności i regularności
miar dodatnich $\mu_x$ na zwartym zbiorze metrycznym $K$.
Jednoznaczność w twierdzeniu Riesza--Markowa daje
równość miar, a więc $E(B)=F(B)$ dla każdego zbioru
borelowskiego $B$.
\end{proof}

\begin{uwaga}[Zakres wykorzystanego twierdzenia miarowego]
Dowód twierdzenia spektralnego sprowadza jego część miarową
do \emph{twierdzenia Riesza--Markowa}: każdy dodatni ciągły
funkcjonał na $C(K)$ dla zwartej przestrzeni Hausdorffa $K$
ma jednoznaczną reprezentację przez skończoną regularną miarę
borelowską. Jest to rezultat skalarnej teorii miary, a nie
wniosek z samego twierdzenia Hahna--Banacha. W pozostałych
krokach dowodu konstrukcja miary projekcyjnej i jej jednoznaczność
zostały przeprowadzone jawnie. Nie używaliśmy założenia
ośrodkowości przestrzeni Hilberta.
\end{uwaga}

\subsection{Konsekwencje miary spektralnej}

\begin{twierdzenie}[Nośnik, widmo i wartości własne]
\label{thm:spectral-support32}
Niech $T$ i $E$ będą jak w twierdzeniu~\ref{thm:spectral-normal32}.
Wtedy $\sigma(T)$ jest najmniejszym domkniętym zbiorem
$F\subseteq K$ takim, że $E(F)=I_H$.
Dla każdego $\lambda\in\C$:
$\lambda$ jest wartością własną $T$ wtedy i tylko wtedy,
gdy $E(\{\lambda\}\cap K)\ne0$, a wówczas
$\ker(T-\lambda I_H)=E(\{\lambda\}\cap K)H$.
Dla $\lambda\in\sigma(T)$ i $\varepsilon>0$ istnieje
$\|x\|=1$ takie, że $\|(T-\lambda I_H)x\|<\varepsilon$.
\end{twierdzenie}
\begin{proof}
Jeśli $U\subseteq\C$ jest otwarty i $U\cap K\ne\varnothing$,
wybieramy niezerową funkcję $f\in C(K)$ o nośniku
zawartym w $U\cap K$ (z odpowiednią interpretacją
nośnika względnego w $K$).
Izometria rachunku ciągłego daje $f(T)\ne0$.
Gdyby $E(U\cap K)=0$, tożsamość
$f(T)=f(T)E(U\cap K)$ dawałaby sprzeczność.
Każdy względnie otwarty niepusty podzbiór $K$
ma zatem niezerową projekcję spektralną. Stąd
$K$ jest nośnikiem $E$.

Jeśli $Tx=\lambda x$, to
$0=\|(T-\lambda I_H)x\|^2
=\int_K|z-\lambda|^2\,\mathrm d\mu_x(z)$,
więc dla każdego $n\ge1$ zbiór
$\{z\in K:|z-\lambda|\ge1/n\}$ ma $\mu_x$-miarę zero.
Ich suma to $K\setminus\{\lambda\}$, stąd
$\|(I_H-E(\{\lambda\}\cap K))x\|^2=0$,
czyli $x=E(\{\lambda\}\cap K)x$.
Odwrotnie, na obrazie tego rzutu operator $T$ działa
jako mnożenie przez $\lambda$.
Wybierając niezerowy wektor z
$E(K\cap B(\lambda,\varepsilon/2))H$ i normalizując go,
otrzymujemy wymagany wektor przybliżony dzięki wzorowi
na normę z twierdzenia~\ref{thm:spectral-normal32}:
$\|(T-\lambda I_H)x\|\le\varepsilon/2<\varepsilon$.
\end{proof}

\begin{twierdzenie}[Borelowski rachunek funkcyjny i zbieżność]
\label{thm:borel-convergence32}
Niech $T\in\mathcal L(H,H)$ będzie normalny z miarą spektralną $E$
na $K=\sigma(T)$. Dla ograniczonej funkcji borelowskiej
$f:K\to\C$ operator $f(T)$ zależy jedynie od wartości $f$
poza zbiorem borelowskim $N$ z $E(N)=0$, a
\[
 \|f(T)\|=\inf\{M\ge0:E(\{z\in K:|f(z)|>M\})=0\}.
\]
Jeżeli $f_n$ są ograniczonymi funkcjami borelowskimi,
$\sup_n\|f_n\|_\infty<\infty$ oraz $f_n(z)\to f(z)$
dla ograniczonej borelowskiej funkcji $f$ poza zbiorem
$E$-zerowym, to $f_n(T)x\to f(T)x$
w normie $H$ dla każdego $x\in H$.
W szczególności $0\le f\le g$ pociąga $f(T)\le g(T)$
dla rzeczywistych ograniczonych funkcji borelowskich $f,g$.
\end{twierdzenie}
\begin{proof}
Jeśli $f=g$ poza zbiorem $N$ z $E(N)=0$, to
$\mu_x(N)=0$ dla wszystkich $x$, a wzór na normę
zastosowany do $f-g$ daje $f(T)=g(T)$.
W twierdzeniu o zbieżności zakładamy, że $f$ jest
ograniczona i borelowska; ewentualnie można przyjąć
ją równą zero na wyjątkowym zbiorze borelowskim.
Dla $x\in H$ wzór na normę daje
$\|(f_n(T)-f(T))x\|^2
=\int_K|f_n(z)-f(z)|^2\,\mathrm d\mu_x(z)$.
Zbieżność zmajoryzowana dowodzi tezy o zbieżności.
Jeżeli $E(\{|f|>M\})=0$, to całki normowe dają
$\|f(T)\|\le M$. Odwrotnie, gdy $E(\{|f|>M\})\ne0$,
możemy wybrać wektor jednostkowy z obrazu tego rzutu;
wtedy $\mu_x$ jest skupiona na $\{|f|>M\}$ i ma masę
$1$. Całka z $|f|^2-M^2$ jest ściśle dodatnia:
zbiór, na którym funkcja ta jest dodatnia, jest sumą
zbiorów, na których jest nie mniejsza niż $1/n$,
a przynajmniej jeden ma dodatnią miarę. Stąd
$\|f(T)x\|>M$. Przechodząc do infimum,
otrzymujemy wzór na normę.
Pozytywność wynika z
$\langle f(T)x,x\rangle=\int_Kf\,\mathrm d\mu_x$.
\end{proof}

\begin{wniosek}[Operatory samosprzężone, dodatnie i unitarne]
\label{cor:spectral-types32}
Niech $T\in\mathcal L(H,H)$ będzie normalny.
Wówczas $T=T^*$ wtedy i tylko wtedy, gdy
$\sigma(T)\subseteq\R$; jest dodatni wtedy i tylko wtedy, gdy
$\sigma(T)\subseteq[0,\infty)$, a unitarny wtedy i tylko wtedy,
gdy $\sigma(T)\subseteq\{z\in\C:|z|=1\}$.
Dla samosprzężonego $T$ projekcje
$E(( -\infty,t]\cap\sigma(T))$, $t\in\R$, tworzą
monotoniczną rodzinę rzutów opisującą rozkład spektralny.
\end{wniosek}
\begin{proof}
Z izometrii rachunku ciągłego wynika, że $T-T^*=0$
dokładnie wtedy, gdy $z-\overline z=0$ na $K$.
Podobnie dla normalnego $T$ oba iloczyny $T^*T$ i $TT^*$
odpowiadają funkcji $|z|^2$, więc unitarność jest
równoważna warunkowi $|z|=1$ na $K$.
Jeśli $K\subseteq[0,\infty)$, to całka z funkcji
$z$ względem każdej $\mu_x$ jest nieujemna, zatem $T\ge0$.
Odwrotnie, gdy $T\ge0$, jest samosprzężony, więc $K\subseteq\R$.
Jeśli $\lambda\in K$ i $\lambda<0$, wybieramy względnie
otwarte otoczenie $V$ punktu $\lambda$ w $K$ zawarte
w $(-\infty,\lambda/2)$. Z Twierdzenia~\ref{thm:spectral-support32}
mamy $E(V)\ne0$. Dla jednostkowego $x\in E(V)H$ miara
$\mu_x$ jest skupiona na $V$ i ma masę $1$, więc
$\langle Tx,x\rangle\le\lambda/2<0$, sprzeczność.
Dla operatora samosprzężonego monotoniczność wynika z
$E(B)E(C)=E(B)$ dla $B\subseteq C$.
\end{proof}

\begin{przyklad}[Widmo ciągłe operatora mnożenia]
\label{ex:mult-spectrum32}
W zespolonej przestrzeni Hilberta $H=L^2([0,1],\lambda_1)$
niech $(Tf)(t)=tf(t)$. Operator jest samosprzężony,
bo mnożnik jest rzeczywisty. Dla $\lambda\notin[0,1]$
odwrotność $\lambda I-T$ jest mnożeniem przez
ograniczoną funkcję $(\lambda-t)^{-1}$.
Dla $\lambda\in[0,1]$ i $\varepsilon>0$ zbiór
$A=[0,1]\cap(\lambda-\varepsilon/2,\lambda+\varepsilon/2)$
ma dodatnią miarę. Funkcja
$f=\chi_A/\sqrt{\lambda_1(A)}$ ma normę $1$ i
$\|(\lambda I-T)f\|\le\varepsilon/2$, co wyklucza
istnienie ograniczonej odwrotności.
Zatem $\sigma(T)=[0,1]$. Miara spektralna ma postać
$(E(B)f)(t)=\chi_{B\cap[0,1]}(t)f(t)$.
Rzuty mnożenia przez funkcje charakterystyczne
spełniają aksjomaty miary spektralnej, w tym silną
addytywność z twierdzenia o zbieżności zmajoryzowanej.
Całka funkcji $z$ jest mnożeniem przez $t$ przez
jednostajną aproksymację funkcjami prostymi, więc
jednoznaczność miary spektralnej uzasadnia podany wzór.
Każda pojedyncza wartość $\lambda\in[0,1]$ ma miarę Lebesgue'a
równą $0$, więc $E(\{\lambda\})=0$ i $T$ nie ma wartości
własnych. Dowodzi to, że nawet dla operatora samosprzężonego
pełny rozkład na wektory własne nie musi istnieć.
\end{przyklad}

\begin{przyklad}[Nieośrodkowa przestrzeń Hilberta]
\label{ex:nonsep-spectrum32}
Niech $I$ będzie dowolnym niepustym zbiorem i $H=\ell^2(I)$.
Dla ograniczonej funkcji $a:I\to\C$ definiujemy
$Te_i=a(i)e_i$ na kanonicznej bazie ortonormalnej.
Współrzędnie $(Tx)_i=a(i)x_i$, więc
$\|Tx\|\le\sup_i|a(i)|\|x\|$ i $T^*$ jest mnożeniem
przez $\overline a$. Stąd normalność.
Poza $\overline{a(I)}$ odwrotność $\lambda I-T$ jest
mnożeniem przez ograniczoną funkcję $i\mapsto(\lambda-a(i))^{-1}$.
Dla $\lambda\in\overline{a(I)}$ istnieją indeksy $i_n$
z $|a(i_n)-\lambda|<1/n$, więc
$\|(\lambda I-T)e_{i_n}\|<1/n$ wyklucza ograniczoną
odwrotność. Zatem $\sigma(T)=\overline{a(I)}$.
Rzut $E(B)$ ma współrzędne $\chi_{a^{-1}(B)}(i)x_i$;
mnożenie tych funkcji charakterystycznych daje aksjomaty
algebraiczne, a zbieżność zmajoryzowana po nośniku
współrzędnych $x$ daje silną addytywność.
Całka funkcji $z$ daje $T$ przez aproksymację prostą,
więc z jednoznaczności miary jest to rzut na domkniętą
rozpiętość $\{e_i:a(i)\in B\}$. Każdy wektor ma co najwyżej przeliczalnie wiele
niezerowych współczynników: dla każdego $n\ge1$ zbiór
$\{i:|x_i|\ge1/n\}$ jest skończony, a suma tych zbiorów
obejmuje wszystkie niezerowe współczynniki.
Sam zbiór indeksów $I$ może być nieprzeliczalny.
\end{przyklad}

\subsection{Rodzina spektralna i rozkład polarny}

\begin{twierdzenie}[Rodzina spektralna operatora samosprzężonego]
\label{thm:spectral-family32}
Niech $T\in\mathcal L(H,H)$ będzie samosprzężony, $E$ jego
miarą spektralną na $K=\sigma(T)\subseteq\R$, a
$P_t=E(K\cap(-\infty,t])$ dla $t\in\R$.
Wtedy $P_sP_t=P_{\min(s,t)}$, $P_s\le P_t$ dla $s\le t$,
$\lim_{t\to-\infty}P_tx=0$ i
$\lim_{t\to+\infty}P_tx=x$ dla wszystkich $x\in H$.
Dla każdego $t\in\R$ i każdego ciągu $t_n\downarrow t$
zachodzi $P_{t_n}x\to P_tx$ w normie; dla $t_n<t$
i $t_n\uparrow t$ granicą jest $E(K\cap(-\infty,t))x$.
Ponadto
\[
 \langle Tx,y\rangle=\int_{K}t\,\mathrm d\mu_{x,y}(t)
 \quad(x,y\in H),
\]
a różnica $P_t-E(K\cap(-\infty,t))=E(K\cap\{t\})$
jest rzutem na $\ker(T-tI_H)$; jej rząd (być może
nieskończony) jest krotnością geometryczną wartości własnej $t$.
\end{twierdzenie}
\begin{proof}
Monotoniczność i równość iloczynów wynikają z własności
$E(B\cap C)=E(B)E(C)$. Ponieważ $K$ jest zwarty,
dla dostatecznie małego $t$ mamy $P_t=0$, a dla dostatecznie
dużego $t$ mamy $P_t=I_H$.
Jeżeli $t_n\downarrow t$, to zbiory $K\cap(-\infty,t_n]$
maleją do $K\cap(-\infty,t]$. Dla każdego $x\in H$
\[
 \|(P_{t_n}-P_t)x\|^2
 =\mu_x(K\cap(t,t_n])\longrightarrow0
\]
na mocy ciągłości miary skończonej z góry.
Analogicznie dla $t_n<t$ i $t_n\uparrow t$ zbiory
$K\cap(-\infty,t_n]$ rosną do $K\cap(-\infty,t)$.
Wzór na parowanie jest szczególnym przypadkiem
rachunku spektralnego, a ostatnia równość wynika
z addytywności miary i z twierdzenia~\ref{thm:spectral-support32}.
\end{proof}

\begin{twierdzenie}[Rozkład polarny operatora ograniczonego]
\label{thm:polar-decomp32}
Niech $H$ będzie zespoloną przestrzenią Hilberta i
$T\in\mathcal L(H,H)$ dowolnym operatorem (niekoniecznie normalnym).
Istnieje dokładnie jedna \emph{częściowa izometria}
$U\in\mathcal L(H,H)$, której przestrzenią początkową jest
$(\ker T)^\perp$ i przestrzenią końcową
$\overline{\operatorname{Ran}T}$, taka że
$T=U|T|$, gdzie $|T|=(T^*T)^{1/2}$ jest jedynym dodatnim
pierwiastkiem z $T^*T$. Co więcej, $\ker U=\ker T$
i $U^*U$ jest rzutem ortogonalnym na $(\ker T)^\perp$.
Dla operatora normalnego $T$ czynniki $U$ i $|T|$ komutują.
\end{twierdzenie}
\begin{proof}
Ponieważ $T^*T\ge0$, ciągły rachunek funkcyjny
dostarcza dodatniego operatora $|T|$ spełniającego
$|T|^2=T^*T$. Uzasadnijmy jednoznaczność dodatniego
pierwiastka dowolnego dodatniego operatora $A$.
Jeśli $B\ge0$ i $B^2=A$, wybieramy wielomiany $p_n$
zbiegające jednostajnie do $\sqrt{\cdot}$ na
$[0,\|A\|]$ (dla $A=0$ mamy $\|B\|^2=\|B^2\|=0$).
Ponieważ $\sigma(B)\subseteq[0,\|B\|]$ oraz
$\|B\|^2=\|A\|$, rachunek ciągły dla $B$ daje
$p_n(B^2)\to B$: funkcje $p_n(t^2)$ zbiegają
jednostajnie do $t$ na $\sigma(B)$.
Rachunek dla $A$ daje jednocześnie $p_n(A)\to A^{1/2}$,
więc $B=A^{1/2}$. Wówczas
$\||T|x\|^2=\langle T^*Tx,x\rangle=\|Tx\|^2$, więc
$\ker|T|=\ker T$ oraz
$\overline{\operatorname{Ran}|T|}=(\ker T)^\perp$.
W ostatniej równości użyliśmy
$(\operatorname{Ran}|T|)^\perp=\ker|T|^*=\ker|T|$.
Na $\operatorname{Ran}|T|$ definiujemy $U_0(|T|x)=Tx$.
Jest to dobrze określona izometria, którą można przedłużyć
przez ciągłość na $(\ker T)^\perp$.
Na $\ker T$ kładziemy $U=0$. Otrzymujemy częściową
izometrię o wskazanych przestrzeniach, normie co najwyżej $1$
i z własnością $U|T|=T$. Jej obraz na przestrzeni
początkowej jest domknięty (izometria z przestrzeni
zupełnej) i równy $\overline{\operatorname{Ran}T}$.
Dla rozkładu $x=x_1+x_0$ z $x_1\perp\ker T$,
$x_0\in\ker T$, mamy $Ux=Ux_1$, zatem $\ker U=\ker T$.
Dla $y=y_1+y_0$ analogicznie
$\langle U^*Ux,y\rangle=\langle x_1,y_1\rangle
=\langle x_1,y\rangle$, więc $U^*U$ jest wskazanym rzutem.
Każda inna częściowa izometria spełniająca warunki tezy
musi zgadzać się z $U$ na gęstym zbiorze
$\operatorname{Ran}|T|$ w przestrzeni początkowej oraz znikać
na $\ker T$, co dowodzi jednoznaczności.
Jeżeli $T$ jest normalny, w modelu spektralnym
$|T|$ odpowiada mnożeniu przez $|z|$, a $U$ przez
funkcję $u(z)=z/|z|$ dla $z\ne0$ i $u(0)=0$:
operator $u(T)$ znika na $\ker T$, jest izometrią na
jej dopełnieniu z wzoru na normę i spełnia
$u(T)|T|=T$; jednoznaczność identyfikuje go z $U$.
Ponieważ funkcje $u(z)$ i $|z|$ komutują punktowo,
operatory $U$ i $|T|$ również komutują.
\end{proof}

\begin{przyklad}[Rozkład polarny macierzy]
Dla $T=\left(\begin{smallmatrix}0&2\\0&0\end{smallmatrix}\right)$
na $\C^2$ mamy
$|T|=\operatorname{diag}(0,2)$ oraz
$U=\left(\begin{smallmatrix}0&1\\0&0\end{smallmatrix}\right)$.
Zatem $T=U|T|$, a $U$ jest częściową izometrią;
nie jest operatorem unitarnym, gdyż $U^*U\ne I$.
\end{przyklad}

\subsection{Operator normalny zwarty jako przypadek dyskretny}

\begin{twierdzenie}[Twierdzenie spektralne dla zwartego operatora normalnego]
\label{thm:compact-normal32}
Niech $H$ będzie niezerową zespoloną przestrzenią Hilberta,
a $T\in\mathcal L(H,H)$ operatorem zwartym i normalnym.
Niezerowe wartości widma tworzą zbiór co najwyżej przeliczalny,
każda z nich jest wartością własną o skończonej krotności,
a jedynym możliwym punktem skupienia tego zbioru jest $0$.
Przestrzeń $H$ jest sumą ortogonalną przestrzeni własnych $T$,
włącznie z $\ker T$. Dla różnych niezerowych wartości
własnych $\lambda$ oznaczmy przez $P_\lambda$ rzut na
$\ker(T-\lambda I_H)$. Wtedy
$T=\sum_{\lambda\in\sigma(T)\setminus\{0\}}
\lambda P_\lambda$, przy czym szereg jest zbieżny w normie
operatorowej (sumy rozumiemy po skończonych podzbiorach indeksów).
\end{twierdzenie}
\begin{proof}
Dla nieskończenie wymiarowego $H$ stosujemy
Twierdzenie~\ref{thm:riesz-schauder30}: niezerowe punkty
widma są izolowanymi wartościami własnymi skończonej
krotności i mają jedyny możliwy punkt skupienia w zerze.
Zbiory punktów o module co najmniej $1/n$ są skończone,
więc ich suma jest co najwyżej przeliczalna.
Dla skończenie wymiarowego $H$ widmo jest skończonym
zbiorem pierwiastków wielomianu charakterystycznego,
a wszystkie przestrzenie własne mają skończony wymiar.
Z twierdzenia~\ref{thm:spectral-support32} każda niepusta
względnie otwarta część widma ma niezerową projekcję spektralną.
Poza zerem widmo jest zatem sumą pojedynczych punktów,
a odpowiadające im projekcje są ortogonalne.
Dopełnienie ortogonalne ich obrazów jest obrazem
$E(K\cap\{0\})$ i równa się $\ker T$.
Dla skończonego zbioru niezerowych wartości własnych $F$
wzór z Twierdzenia~\ref{thm:borel-convergence32} daje
\[
\left\|T-\sum_{\lambda\in F}\lambda P_\lambda\right\|
=\sup\bigl(\{0\}\cup\{|\lambda|:\lambda\in K\setminus(F\cup\{0\})\}\bigr).
\]
Istotnie, każda pozostająca niezerowa wartość własna
ma niezerową projekcję, a reszta funkcji spektralnej
jest zerowa na $F$ i w zerze. Wzór obejmuje też
przypadek, gdy nie pozostaje żadna niezerowa wartość.
Ta wielkość dąży do zera, gdy $F$ obejmuje wszystkie
wartości własne o module większym od ustalonego $\varepsilon>0$.
W rezultacie otrzymujemy zbieżność w normie operatorowej.
\end{proof}

\begin{uwaga}[Granica twierdzenia]
Dla operatorów nienormalnych nie można w ogólności znaleźć
miary spektralnej projekcyjnej reprezentującej operator.
Na przykład niezerowa macierz
$\left(\begin{smallmatrix}0&1\\0&0\end{smallmatrix}\right)$
ma widmo $\{0\}$, lecz nie jest zerowa; gdyby miała reprezentację
$T=\int z\,\mathrm dE(z)$ z miarą spektralną na tym widmie,
byłaby równa $0$. Rozszerzenie rachunku spektralnego na
operatory nieograniczone opracujemy w \hyperref[sec:unbounded33]{rozdziale~33}.
\end{uwaga}
